% Source: https://artofproblemsolving.com/community/c3h1294839_infinite_series
% Stable identity: aops-topic-1294839-q01
The argument treats $S$ as a number, but the series has no (finite) sum. The partial sums are
\[
S_n=1+3+\cdots+3^n=\dfrac{3^{n+1}-1}{2},
\]
which grow without bound ($S_n\ge3^n$). An infinite sum is \emph{defined} as the limit of the partial sums, and here there is no limit, so ``$S$'' is not a real number. The rules ``$3S$ is the series multiplied term by term'' and ``$S-3S$ can be computed term by term'' are only valid for convergent series; applied to a divergent one they produce nonsense (the same steps applied to $T=1+2+4+\cdots$ would give $T=-1$, a negative ``sum'' of positive numbers). So the error is the very first step, assuming the sum $S$ exists.

Done correctly with finite sums, the same idea works: $3S_n-S_n=3^{n+1}-1$, which gives the formula for $S_n$ above. For a ratio $r$ with $|r|<1$ the partial sums $\dfrac{1-r^{n+1}}{1-r}$ do converge, and then the trick legitimately gives $1+r+r^2+\cdots=\dfrac{1}{1-r}$; for $r=3$ it does not.
